\begin{proof}
Let  

\[
f(t)=\sum_{i=1}^{n}\bigl(a_i-t\,b_i\bigr)^{2},\qquad t\in\mathbb{R}.
\]

Each summand is a square, hence \(f(t)\ge0\) for all \(t\).  

Expanding \(f(t)\) using \((a-tb)^{2}=a^{2}-2tab+t^{2}b^{2}\) gives  

\[
\begin{aligned}
f(t)
&=\sum_{i=1}^{n}\Bigl(a_i^{2}-2t a_i b_i+t^{2}b_i^{2}\Bigr)\\[2mm]
&=\Bigl(\sum_{i=1}^{n}b_i^{2}\Bigr)t^{2}
   -2\Bigl(\sum_{i=1}^{n}a_i b_i\Bigr)t
   +\sum_{i=1}^{n}a_i^{2}.
\end{aligned}
\]

Set  

\[
A:=\sum_{i=1}^{n}b_i^{2},\qquad 
B:=-2\sum_{i=1}^{n}a_i b_i,\qquad 
C:=\sum_{i=1}^{n}a_i^{2},
\]

so that \(f(t)=At^{2}+Bt+C\).

Since \(f(t)\ge0\) for every real \(t\), the quadratic polynomial must have non‑positive discriminant:

\[
B^{2}-4AC\le0. \tag{1}
\]

Substituting the expressions for \(A,B,C\) into (1) yields  

\[
\bigl(-2\sum_{i=1}^{n}a_i b_i\bigr)^{2}
-4\Bigl(\sum_{i=1}^{n}b_i^{2}\Bigr)\Bigl(\sum_{i=1}^{n}a_i^{2}\Bigr)\le0.
\]

Dividing by \(4>0\) and simplifying gives  

\[
\bigl(\sum_{i=1}^{n}a_i b_i\bigr)^{2}
\le
\Bigl(\sum_{i=1}^{n}a_i^{2}\Bigr)
\Bigl(\sum_{i=1}^{n}b_i^{2}\Bigr),
\]

which is the Cauchy–Schwarz inequality.

\medskip
\noindent\textbf{Equality condition.}  
Equality in (1) occurs iff the discriminant is zero, i.e. the quadratic has a double root. In that case  

\[
f(t)=A\bigl(t-t_{0}\bigr)^{2},\qquad 
t_{0}=\frac{\sum_{i=1}^{n}a_i b_i}{\sum_{i=1}^{n}b_i^{2}}.
\]

Because \(f(t)=\sum_{i=1}^{n}(a_i-t b_i)^{2}\), the equality \(f(t_{0})=0\) forces each term to vanish:

\[
a_i=t_{0}\,b_i\quad\text{for all }i.
\]

Conversely, if there exists \(\lambda\in\mathbb{R}\) with \(a_i=\lambda b_i\) for every \(i\), then every summand \((a_i-\lambda b_i)^{2}\) is zero, so \(f(\lambda)=0\) and the discriminant is zero. Hence equality holds precisely when the vectors \((a_{1},\dots ,a_{n})\) and \((b_{1},\dots ,b_{n})\) are linearly dependent. \(\square\)
\end{proof}